\documentclass[b5paper,twoside,10pt]{book}
\usepackage{natbib}
\usepackage[utf8]{inputenc}
\usepackage[T1]{fontenc}
\usepackage{lmodern}
\usepackage[spanish]{babel}
\usepackage{xspace}
\usepackage{amsmath}
\usepackage{framed}
\usepackage{type1cm}
\usepackage{textcomp} %para \textcelsius
\usepackage[indent,small,it,bf]{caption}
%formato de los captions (pies de tablas y figuras): ``using imported graphics in latex (pág 50).
\bibliographystyle{estilo}
%\usepackage[nottoc,notlot,notlof]{tocbibind}
% para que la bibliografía salga en el índice y no lo haga el lot y el lof.
\usepackage{multicol}
% para juntar columnas en las tablas.
\usepackage{marvosym,amssymb,url}
% para insertar símbolos como \Male y \Female.
\usepackage{setspace,paralist,mdwlist}
%personalizar listas
\usepackage{graphicx}}
%paquetes para insertar imágenes.
\usepackage{afterpage}
%para colocar un float en la pagina siguiente
\usepackage[table]{xcolor}
%permite poner colores en las tablas
\usepackage{multirow}
%para juntar filas en las tablas
\usepackage{rotating}
%para rotar texto

%\usepackage{texshade}
% alineaciones de secuencias etc.. requiere paso por postscript
\usepackage{longtable}
%\usepackage[paperwidth=17.5cm,paperheight=25cm,top=3cm,bottom=2.75cm,inner=3cm,outer=2.5cm,dvips=false,pdftex=false,vtex=false]{geometry}
\usepackage[]{geometry}
\geometry{
%paperwidth=17.5cm,
%paperheight=25cm,
top=3cm,
bottom=3.25cm,
inner=3.25cm,
outer=2.5cm
}

%\usepackage[cam,a4,center,dvips]{crop} %ejes
\usepackage[a4,center,dvips]{crop} %centrada
%\usepackage{crop} %esquina

\usepackage{hyperref}
\hypersetup{colorlinks=true,linkcolor=red,citecolor=black,urlcolor=red,bookmarks=false}
%\hypersetup{colorlinks=true,linkcolor=black,citecolor=black,urlcolor=black,bookmarks=false}
%\usepackage{tocloft}
%para personalizar la toc.

\usepackage{epsfig}
%para que latex pueda rotular las figuras. Requiere paso por postscript. Uso de xfig
\usepackage{texshade}
%(para crear los alineamientos de secuencias)
\usepackage{hhline}

\usepackage{fancyvrb}

\date{8 de Julio de 2011}

\title{Aquí va el título de la tesis}

\author{Nombre y Apellidos}
%\setlength{\parskip}{.7ex}
\setlength\parskip{.5\baselineskip
	plus .1\baselineskip
		minus .4\baselineskip
		}
% sepaqración entre párrafos.

\setcounter{tocdepth}{2}
% profundidad en la tabla de contenidos.

% \setcounter{secnumdepth}{5}
% numera todos los epígrafes hasta nivel \subparagraph en el cuerpo del documento


\makeatletter
\def\@makechapterhead#1{%
  \vspace*{1\p@}%
  {\parindent \z@ \raggedright \normalfont
    \ifnum \c@secnumdepth >\m@ne
%        \huge\bfseries \@chapapp\space \thechapter
        \par\nobreak
        \vskip 20\p@
    \fi
    \interlinepenalty\@M
    \Huge \bfseries \thechapter.\space#1\par\nobreak
    \vskip 40\p@
}}
\makeatother
% Cambia el formato de los títulos de los capítulos.

%serie de tamaños de fuente: \tiny, \scriptsize, \footnotesize, \small, \normalsize, \large, \Large, \LARGE, \huge, \Huge

\addtolength{\headheight}{12pt}
%%%%%%%%%%%%%%%%%%%%%%%%%%%%%%%%%%%%%%%%%%%
%                                         %
%   Para encabezados y pies de página     %
%                                         %
%%%%%%%%%%%%%%%%%%%%%%%%%%%%%%%%%%%%%%%%%%%
\usepackage{fancyhdr}
%\pagestyle{plain} 
\renewcommand{\footrulewidth}{0pt}
\fancyhf{}

%\usepackage{glosstex}
%genera el glosario.

%\sloppy
%suaviza las rupturas de línea.

%The \flushbottom declaration makes all text pages the same height,adding extra vertical space when necessary to fill out the page. This is the standard for the book style.

% The \raggedbottom declaration makes all pages the height of the text on that page. No extra vertical space is added.

\pretolerance=1000
\tolerance=2000
%(para que no corte muchas palabras)

%\usepackage{titlesec}
%\titleformat*{\section}{\bfseries \Large}
%\titleformat*{\subsection}{\bfseries \large}


\newcommand{\um}{\ensuremath{\mu m}}

%\newcommand{\subsubsubsection}[1]{\paragraph{\vspace{1ex}#1\vspace{1ex}}}
%\newcommand{\apartado}[1]{\vspace{1ex}\noindent\textbf{\emph{#1}}\vspace{1ex}}
%\newcommand{\titulillo}[1]{\vspace{1ex}\textbf{\Large #1}\vspace{1ex}}
\newcommand{\GC}{\textcelsius\xspace}
\newcommand{\ul}{\ $\mu$l\xspace}
\newcommand{\ug}{\ $\mu$g\xspace}
\newcommand{\infi}{$\infty$}
\newcommand{\uM}{$\mu$M}

\newcommand{\gris}[2]{\multicolumn{1}{>{\columncolor[gray]{0.9}}#1}{#2}}
%\newcommand{\gris}[2]{\multicolumn{1}{#1}{#2}}

%\newcommand{\epigrafe}[1]{\textbf{\emph{#1}}}

%\setlength{\extrarowheight}{-.5cm}


%para poner las figuras con solo \figura{nombre archivo sin extensión}{tamaño .8}{titulo corto xa el índice}{pie figura}
\usepackage{float}
\newcommand{\figura}[4]{%
   \begin{figure}[htbp]
    \begin{center}
     \includegraphics[width=#2\textwidth]{figuras/#1}
    \end{center}
     \caption[#3]{#4}
     \label{#1}
    \end{figure}}


%\usepackage[none]{hyphenat}
%		\hyphenation{mRNA mRNAs re-gu-la-do e-sen-ci-al}

\hyphenation{mRNA miRNA miRNAs mRNAs}


\newcommand{\titulo}{Este es el nuevo título de mi Tesis}
\newcommand{\fecha}{22 julio 2016}

%%%%%%%%%%%%%%%%%%%%%%%%%%%%%%%%%%%%%%%%%%%%%%%%%%%
%                                                 %
%                BEGIN DOCUMENT                   %
%                                                 %
%%%%%%%%%%%%%%%%%%%%%%%%%%%%%%%%%%%%%%%%%%%%%%%%%%%
\begin{document}
%\maketitle


\frontmatter

\renewcommand{\labelitemi}{$\bullet$}
\renewcommand{\labelitemii}{---}
\renewcommand{\labelitemiii}{--}


\pagestyle{empty}
\begin{center}
\LARGE
\titulo
\end{center}
%\vspace{20ex}

%\vfill

\vspace{\stretch{1}}

\begin{figure}[h]
\centering
\includegraphics[width=.6\textwidth]{escudo_ugr_portada}
\end{figure}

%\vspace{25ex}
%\vfill

\vspace{\stretch{2}}

\begin{center}
\large
Nombre Apellidos
\end{center}
\begin{center}
\large
Tesis Doctoral

\fecha
\end{center}
%\newpage
\cleardoublepage
\thispagestyle{empty}
\begin{minipage}[t]{.23\textwidth}
\centering
\includegraphics[width=.8\textwidth]{logo_ib}
Instituto de Biotecnolog\'ia 
\end{minipage}
\hfill
\begin{minipage}[t]{.23\textwidth}
\centering
\includegraphics[width=.7\textwidth]{logo_ciencias}
Facultad de Ciencias
\end{minipage}
\hfill
\begin{minipage}[t]{.23\textwidth}
\centering
\includegraphics[width=.83\textwidth]{escudo_ugr_portada}
 Universidad de Granada
\end{minipage}


\begin{minipage}[b]{.23\textwidth}
\centering
\includegraphics[width=.7\textwidth]{biologia}
Estudios de Biolog\'ia
\end{minipage}
\hfill
\begin{minipage}[b]{.23\textwidth}
\centering
\includegraphics[width=.7\textwidth]{genetica}
Departamento de Gen\'etica
\end{minipage}
\hfill
\begin{minipage}[b]{.23\textwidth}
\centering
\includegraphics[width=.7\textwidth]{cibm}
Centro de Investigaci\'on Biom\'edica
\end{minipage}
\vfill
\begin{center}
\LARGE
\titulo

\end{center}
\vfill

\begin{minipage}[t]{.4\textwidth}
Dirigida por los Doctores:

\vspace{15ex}

Dr. Fulanito de tal

\vspace{15ex}

Dr. Menganito de cual

\end{minipage}
\hfill
\begin{minipage}[t]{.5\textwidth}
Memoria de Tesis Doctoral presentada por la Licenciada Nombre Apellidos  para optar al grado de ``Doctor por la Universidad de Granada''.

\vspace{29.7ex}

\centering 
Fulanita de tal

\fecha

\end{minipage}


%\newpage
\cleardoublepage
\thispagestyle{empty}

\null\vspace{\stretch{5}}
El presente trabajo se ha desarrollado en los laboratorios del Departamento de Gen\'etica y del Instituto de Biotecnolog\'ia en la Facultad de Ciencias y en el Centro de Investigaci\'on Biom\'edica (CIBM) de la Universidad de Granada. La investigaci\'on realizada ha sido financiada por la Junta de Andalucía, mediante la concesi\'on de una beca predoctoral. Este trabajo ha sido financiado por la Junta de Andalucía, a través del Programa de Ayudas a Grupos de Investigación (Grupo PAI BIO-109) y de un Proyecto de Excelencia (CVI-2057).
\vspace{\stretch{1}}
 
\cleardoublepage
\section*{Agradecimientos}

Lorem ipsum dolor sit amet, consectetur adipiscing elit. Proin sed erat porttitor, faucibus dolor sit amet, tincidunt elit. Aenean eget est mi. Quisque ultrices, felis consectetur ultricies efficitur, tellus enim suscipit est, vitae scelerisque augue erat quis quam. Vivamus cursus, lectus gravida ornare varius, elit risus elementum libero, eget facilisis augue sem non orci. Aliquam vestibulum felis vitae massa iaculis, eget lobortis arcu gravida. Sed urna nulla, sagittis a diam sit amet, suscipit vulputate ex. Praesent hendrerit ullamcorper odio, eu eleifend velit tempor in. Maecenas eget ultrices risus, nec tempus felis. Fusce tristique diam eu ligula vestibulum sagittis. Suspendisse viverra et nunc sit amet commodo. Vivamus eu accumsan arcu. Pellentesque laoreet convallis elit at vehicula. Etiam lacinia risus eu erat imperdiet lacinia eu vitae augue. Suspendisse dictum urna non augue ultricies, quis maximus risus hendrerit. Nunc mollis nunc est, sit amet feugiat ex porttitor non.

Nulla vitae sollicitudin eros. Aliquam tempor, ex et tincidunt ultrices, lorem enim lacinia elit, eget malesuada urna orci eget dolor. Nulla elementum, massa in porta egestas, velit nisl maximus lectus, a luctus risus lorem a nisi. Suspendisse dapibus purus quam, nec laoreet mi elementum vel. Cras ut felis in urna gravida ultrices. Interdum et malesuada fames ac ante ipsum primis in faucibus. In ullamcorper pharetra urna sit amet hendrerit. Proin placerat ante quis nisi convallis, quis ultrices elit malesuada. Integer ac tortor gravida, luctus arcu sit amet, semper elit. Fusce nulla diam, interdum et leo auctor, fringilla rhoncus nunc. Duis tempus sed nisi ac fringilla. Maecenas in aliquam massa. Morbi at magna vel elit gravida volutpat.

Pellentesque molestie eros mauris, a luctus nunc vestibulum sed. Phasellus molestie consectetur velit, a tempus orci mattis eu. Maecenas felis ligula, consequat vitae consectetur in, scelerisque eu nibh. Suspendisse posuere dolor eu dui facilisis gravida. Vivamus dictum eleifend lacinia. Vestibulum vehicula luctus hendrerit. Maecenas pellentesque urna blandit augue fermentum vehicula. Morbi sit amet tincidunt tortor. Etiam non fermentum enim. Donec sed euismod sapien, sit amet pretium ex.

Suspendisse tristique, justo ac interdum cursus, enim purus pulvinar eros, ac lobortis lectus enim ac mi. Morbi aliquam consectetur laoreet. Sed congue blandit condimentum. Morbi mollis vel enim ut mollis. Pellentesque ipsum nulla, ultrices quis aliquet non, faucibus quis arcu. Curabitur vulputate cursus ligula, ut ultricies odio gravida vitae. Integer sed scelerisque arcu. Nulla in porta felis, id vehicula elit. Vestibulum augue metus, venenatis eget congue eu, ullamcorper at diam. Ut fermentum diam eros, sed luctus neque dignissim nec. Nullam vel lobortis nulla.

\enlargethispage{2em}

Nulla hendrerit ornare vulputate. Praesent pharetra turpis eget urna maximus laoreet. Ut nibh magna, dapibus in ante id, blandit elementum arcu. Fusce cursus, enim a fringilla faucibus, libero justo ornare elit, quis porta felis sapien quis libero.
 Phasellus imperdiet lectus at justo interdum, quis ultricies mi aliquet. 
%Mauris varius diam mollis est lobortis dignissim. Donec molestie est nisl, vel malesuada quam finibus ut. Pellentesque quis ipsum sed nibh tincidunt sollicitudin. Nam sit amet commodo eros. Morbi quis risus at neque pellentesque congue. Donec at viverra leo, ut volutpat massa. 


\cleardoublepage
\null\vspace{\stretch{2}}
\begin{flushright}
 Dedicado a mis Padres 

y también a mis Hermanos
\end{flushright}
\vspace{\stretch{1}}
\cleardoublepage




\null\vspace{\stretch{2}}
{
\hfill \begin{minipage}{8cm}
\textsl{Siembra un pensamiento y cosecharás un acto,}

\textsl{siembra un acto y cosecharás un hábito,}

\textsl{siembra un hábito y cosecharás un carácter,}

\textsl{siembra un carácter y cosecharas un destino. }

\begin{flushright}
Charles Reade  
\end{flushright}

\end{minipage}
}
\vspace{\stretch{1}}
\cleardoublepage






\tableofcontents
\listoffigures
\listoftables
\clearpage
\section*{Abreviaturas}

\begin{description}
\item[ADN:] Ácido desoxirribonucléico
\item[ARN:] Ácido ribonucléico

\end{description}


%\pagestyle{fancy}

\mainmatter

\pagestyle{fancy}
\lhead{\uppercase{\bfseries{\leftmark}}}
\rhead{\uppercase{\bfseries{\rightmark}}}
\fancyhead[LO]{\bfseries{}}
\fancyhead[RE]{\bfseries{}}
%\fancyfoot[RO]{\thepage}
%\fancyfoot[LE]{\thepage}
\cfoot{\thepage}
\renewcommand{\headrulewidth}{0,2pt}
%estilo de las cabeceras y los pies: L=left, R=right, E=even, O=odd (impares).
\chapter{Abstract}


Lorem ipsum dolor sit amet, consectetur adipiscing elit. Proin sed erat porttitor, faucibus dolor sit amet, tincidunt elit. Aenean eget est mi. Quisque ultrices, felis consectetur ultricies efficitur, tellus enim suscipit est, vitae scelerisque augue erat quis quam. Vivamus cursus, lectus gravida ornare varius, elit risus elementum libero, eget facilisis augue sem non orci. Aliquam vestibulum felis vitae massa iaculis, eget lobortis arcu gravida. Sed urna nulla, sagittis a diam sit amet, suscipit vulputate ex. Praesent hendrerit ullamcorper odio, eu eleifend velit tempor in. Maecenas eget ultrices risus, nec tempus felis. Fusce tristique diam eu ligula vestibulum sagittis. Suspendisse viverra et nunc sit amet commodo. Vivamus eu accumsan arcu. Pellentesque laoreet convallis elit at vehicula. Etiam lacinia risus eu erat imperdiet lacinia eu vitae augue. Suspendisse dictum urna non augue ultricies, quis maximus risus hendrerit. Nunc mollis nunc est, sit amet feugiat ex porttitor non.


Nulla vitae sollicitudin eros. Aliquam tempor, ex et tincidunt ultrices, lorem enim lacinia elit, eget malesuada urna orci eget dolor. Nulla elementum, massa in porta egestas, velit nisl maximus lectus, a luctus risus lorem a nisi. Suspendisse dapibus purus quam, nec laoreet mi elementum vel. Cras ut felis in urna gravida ultrices. Interdum et malesuada fames ac ante ipsum primis in faucibus. In ullamcorper pharetra urna sit amet hendrerit. Proin placerat ante quis nisi convallis, quis ultrices elit malesuada. Integer ac tortor gravida, luctus arcu sit amet, semper elit. Fusce nulla diam, interdum et leo auctor, fringilla rhoncus nunc. Duis tempus sed nisi ac fringilla. Maecenas in aliquam massa. Morbi at magna vel elit gravida volutpat.

\url{http://google.es}

Pellentesque molestie eros mauris, a luctus nunc vestibulum sed. Phasellus molestie consectetur velit, a tempus orci mattis eu. Maecenas felis ligula, consequat vitae consectetur in, scelerisque eu nibh. Suspendisse posuere dolor eu dui facilisis gravida. Vivamus dictum eleifend lacinia. Vestibulum vehicula luctus hendrerit. Maecenas pellentesque urna blandit augue fermentum vehicula. Morbi sit amet tincidunt tortor. Etiam non fermentum enim. Donec sed euismod sapien, sit amet pretium ex.

Suspendisse tristique, justo ac interdum cursus, enim purus pulvinar eros, ac lobortis lectus enim ac mi. Morbi aliquam consectetur laoreet. Sed congue blandit condimentum. Morbi mollis vel enim ut mollis. Pellentesque ipsum nulla, ultrices quis aliquet non, faucibus quis arcu. Curabitur vulputate cursus ligula, ut ultricies odio gravida vitae. Integer sed scelerisque arcu. Nulla in porta felis, id vehicula elit. Vestibulum augue metus, venenatis eget congue eu, ullamcorper at diam. Ut fermentum diam eros, sed luctus neque dignissim nec. Nullam vel lobortis nulla.

Nulla hendrerit ornare vulputate. Praesent pharetra turpis eget urna maximus laoreet. Ut nibh magna, dapibus in ante id, blandit elementum arcu. Fusce cursus, enim a fringilla faucibus, libero justo ornare elit, quis porta felis sapien quis libero. Phasellus imperdiet lectus at justo interdum, quis ultricies mi aliquet. Mauris varius diam mollis est lobortis dignissim. Donec molestie est nisl, vel malesuada quam finibus ut. Pellentesque quis ipsum sed nibh tincidunt sollicitudin. Nam sit amet commodo eros. Morbi quis risus at neque pellentesque congue. Donec at viverra leo, ut volutpat massa. 




Lorem ipsum dolor sit amet, consectetur adipiscing elit. Proin sed erat porttitor, faucibus dolor sit amet, tincidunt elit. Aenean eget est mi. Quisque ultrices, felis consectetur ultricies efficitur, tellus enim suscipit est, vitae scelerisque augue erat quis quam. Vivamus cursus, lectus gravida ornare varius, elit risus elementum libero, eget facilisis augue sem non orci. Aliquam vestibulum felis vitae massa iaculis, eget lobortis arcu gravida. Sed urna nulla, sagittis a diam sit amet, suscipit vulputate ex. Praesent hendrerit ullamcorper odio, eu eleifend velit tempor in. Maecenas eget ultrices risus, nec tempus felis. Fusce tristique diam eu ligula vestibulum sagittis. Suspendisse viverra et nunc sit amet commodo. Vivamus eu accumsan arcu. Pellentesque laoreet convallis elit at vehicula. Etiam lacinia risus eu erat imperdiet lacinia eu vitae augue. Suspendisse dictum urna non augue ultricies, quis maximus risus hendrerit. Nunc mollis nunc est, sit amet feugiat ex porttitor non.

Nulla vitae sollicitudin eros. Aliquam tempor, ex et tincidunt ultrices, lorem enim lacinia elit, eget malesuada urna orci eget dolor. Nulla elementum, massa in porta egestas, velit nisl maximus lectus, a luctus risus lorem a nisi. Suspendisse dapibus purus quam, nec laoreet mi elementum vel. Cras ut felis in urna gravida ultrices. Interdum et malesuada fames ac ante ipsum primis in faucibus. In ullamcorper pharetra urna sit amet hendrerit. Proin placerat ante quis nisi convallis, quis ultrices elit malesuada. Integer ac tortor gravida, luctus arcu sit amet, semper elit. Fusce nulla diam, interdum et leo auctor, fringilla rhoncus nunc. Duis tempus sed nisi ac fringilla. Maecenas in aliquam massa. Morbi at magna vel elit gravida volutpat.

Pellentesque molestie eros mauris, a luctus nunc vestibulum sed. Phasellus molestie consectetur velit, a tempus orci mattis eu. Maecenas felis ligula, consequat vitae consectetur in, scelerisque eu nibh. Suspendisse posuere dolor eu dui facilisis gravida. Vivamus dictum eleifend lacinia. Vestibulum vehicula luctus hendrerit. Maecenas pellentesque urna blandit augue fermentum vehicula. Morbi sit amet tincidunt tortor. Etiam non fermentum enim. Donec sed euismod sapien, sit amet pretium ex.

Suspendisse tristique, justo ac interdum cursus, enim purus pulvinar eros, ac lobortis lectus enim ac mi. Morbi aliquam consectetur laoreet. Sed congue blandit condimentum. Morbi mollis vel enim ut mollis. Pellentesque ipsum nulla, ultrices quis aliquet non, faucibus quis arcu. Curabitur vulputate cursus ligula, ut ultricies odio gravida vitae. Integer sed scelerisque arcu. Nulla in porta felis, id vehicula elit. Vestibulum augue metus, venenatis eget congue eu, ullamcorper at diam. Ut fermentum diam eros, sed luctus neque dignissim nec. Nullam vel lobortis nulla.

Nulla hendrerit ornare vulputate. Praesent pharetra turpis eget urna maximus laoreet. Ut nibh magna, dapibus in ante id, blandit elementum arcu. Fusce cursus, enim a fringilla faucibus, libero justo ornare elit, quis porta felis sapien quis libero. Phasellus imperdiet lectus at justo interdum, quis ultricies mi aliquet. Mauris varius diam mollis est lobortis dignissim. Donec molestie est nisl, vel malesuada quam finibus ut. Pellentesque quis ipsum sed nibh tincidunt sollicitudin. Nam sit amet commodo eros. Morbi quis risus at neque pellentesque congue. Donec at viverra leo, ut volutpat massa. 




\chapter{Introducción y Objetivos}
\input{introduccion/objetivos.tex}

\chapter{Antecedentes}
\input{introduccion/Introduccion.tex}

\chapter{Material y Métodos}
\input{mat-met/materialYmetodos.tex}

\chapter{Resultados}
\input{resultados/resultados.tex}

\chapter{Discusión}
\input{discusion/discusion.tex}

\chapter{Conclusiones}
\input{discusion/conclusiones.tex}


%\addcontentsline{toc}{chapter}{Bibliography}
\begin{singlespace}
	\begin{footnotesize}
		\begin{twocolumn}
\bibliography{biblio}
\end{twocolumn}
\end{footnotesize}
\end{singlespace}

\end{document}
